
\section{Introducción}
\label{sec:intro}

Large Language Models (LLMs)~\cite{brown2020language,ouyang2022training,chatgpt,achiam2023gpt,vicuna2023,touvron2023llama2,dubey2024llama,jiang2024mixtral} han revolucionado el campo de la inteligencia artificial y producido cambios significativos en múltiples áreas. Sobre la base de los LLMs, los Multimodal Large Language Models (MLLMs) han permitido recientemente avances importantes en problemas de image~\cite{alayrac2022flamingo,bai2023qwen,chen2023internvl,dong2024internlm,li2023blip2,zhu2023minigpt,liu2024visual,liu2023improved,huang2024lita,ye2023mplugowl, li2024llava,  li2024llavanext-interleave, llama3.2} y video understanding, que normalmente se reformulan como generación de lenguaje a partir de video~\cite{chen2023videollm, lin2023video, li2023videochat, li2023mvbench, liu2024world, Maaz2023VideoChatGPT, wang2024internvideo2, zhang2023video, rawal2024cinepile}.

Sin embargo, la mayoría de los video MLLMs se han aplicado a entradas de video cortas de apenas unos segundos. En este caso, una estrategia simple pero efectiva consiste en extraer ``tokens'' de los frames individuales del video utilizando un pretrained image encoder~\cite{llama3.2, ataallah2024minigpt4, liu2024visual, li2024llava, wang2024tarsier, li2024llavanext-interleave}. Luego, la concatenación de los tokens extraídos de los frames se envía al LLM para el procesamiento semántico y la generación de lenguaje. Esta estrategia deja de funcionar cuando el video es muy largo, por ejemplo, cuando abarca minutos u horas, debido al costo cuadrático de la operación self-attention necesaria para procesar los tokens visuales extraídos. Además, la concatenación ingenua de tokens a nivel de frame para un video largo saturaría al LLM con una secuencia de características extremadamente larga y redundante. Por ejemplo, el image encoder LLAMA-3.2~\cite{llama3.2} produce entre 1600 y 6400 tokens por imagen. Por lo tanto, codificar solo 128 frames produciría una secuencia de 205K a 820K tokens, lo cual es impracticable de procesar incluso con GPUs modernas.

Para abordar estos problemas, la mayoría de los video MLLMs long-form adoptan técnicas de compresión para reducir la cantidad de tokens visuales. Por ejemplo, los métodos de~\cite{xu2024pllavaparameterfreellava, Maaz2023VideoChatGPT} aplican simple spatial/temporal pooling sobre tokens a nivel de frame para acortar la secuencia enviada al LLM. Sin embargo, la operación de pooling descarta información espaciotemporal importante. Una estrategia alternativa consiste en usar módulos basados en convolución para realizar simultáneamente compresión de secuencia y modelado temporal~\cite{ lin2023video, liu2024kangaroo}. No obstante, los modelos basados en convolución carecen de capacidad de modelado de largo alcance, ya que los kernels convolucionales solo pueden capturar dependencias locales de corto alcance. Trabajos recientes~\cite{zhang2023video} reducen la cantidad de tokens visuales de salida aplicando self-attention con una menor cantidad de query keys~\cite{jaegle2021perceiver}. Sin embargo, estos sistemas intercambian eficiencia por análisis entre frames, algo esencial para el long video understanding.

En este trabajo proponemos {\bf B}idirectional {\bf I}nterleaved {\bf M}amba for {\bf B}etter {\bf A}nswers (\model), que ofrece una alternativa eficiente y efectiva para comprimir videos largos en una secuencia corta de tokens rica en información para video question-answering con LLMs. Nuestro sistema se inspira en el éxito de los state-space models (SSMs)~\cite{gu2020hippo, gu2021combining, gu2022efficiently, gu2022s4d, gupta2022diagonal} para procesar documentos largos en NLP~\cite{gu2023mamba}. Los SSMs poseen de forma inherente capacidad de modelado a largo plazo y requieren costos de cómputo lineales con respecto a la longitud de la secuencia, en lugar del costo cuadrático de self-attention. En particular, Mamba~\cite{gu2023mamba} mejoró aún más los SSMs con un mecanismo de selección que permite al modelo seleccionar información relevante de manera dependiente de la entrada. Diseñamos una arquitectura de video basada en el módulo Mamba para comprimir la secuencia de tokens espaciotemporales en más de un orden de magnitud (por ejemplo, de 102K tokens a 6.4k tokens) mientras conserva dependencias esenciales de largo alcance. El mecanismo selective-scan~\cite{gu2023mamba} de nuestro módulo token selector permite al modelo propagar selectivamente tokens importantes y descartar los redundantes dentro del contenido de video altamente redundante que suele estar presente en videos largos.

Aunque los modelos SSM y Mamba ya se han utilizado en image analysis~\cite{zhu2024vision, pei2024efficientvmamba, pei2024efficientvmamba, chen2024rsmamba}, image-language modeling~\cite{qiao2024vl,zhao2024cobra, huang2024clip, lee2024meteor} y video understanding~\cite{yang2024vivim,li2025videomamba,lu2024videomambapro,chaudhuri2024simba,islam2022long,islam2023efficient}, nuestra arquitectura propuesta se diferencia de estos trabajos previos al introducir varias contribuciones simples pero efectivas. Primero, introducimos un módulo de compresión basado en SSM llamado spatiotemporal token selector, que puede tomar como entrada una gran cantidad de tokens de video espaciotemporales y producir una cantidad significativamente menor de tokens comprimidos que contienen solo información importante mediante el mecanismo selective-scan~\cite{gu2023mamba}. Logramos esto introduciendo una menor cantidad de visual queries y modelándolas conjuntamente con los tokens de video espaciotemporales. Segundo, proponemos un método mejorado para concatenar tokens visuales y espaciotemporales intercalándolos a intervalos iguales. El enfoque estándar de anexar los tokens al final puede introducir sesgo posicional, donde los tokens del final de la secuencia influyen de manera desproporcionada sobre los query tokens. Nuestro posicionamiento interleaved mitiga este sesgo al influir sobre las queries de manera más uniforme a lo largo de toda la secuencia de video. Finalmente, usamos una estrategia bidirectional selective-scan~\cite{li2025videomamba} que resulta más efectiva para capturar estructuras espaciotemporales 2D/3D en video, en comparación con el selective scan estándar, que se adapta mejor al modelado de secuencias 1D.


Experimentamos en varios datasets de long video question answering, incluidos PerceptionTest~\cite{patraucean2024perception}, NExT-QA~\cite{xiao2021next}, EgoSchema~\cite{mangalam2023egoschema}, VNBench~\cite{zhao2024needle}, LongVideoBench~\cite{wu2024longvideobench}, Video-MME~\cite{fu2024video} y MLVU~\cite{zhou2024mlvu}, y mostramos resultados state-of-the-art en todos ellos.
